\theory{Module theory}{What remains true about linear algebra when scalars come from a ring instead of a field?}{A module is an abelian group on which a ring acts like scalars. Modules generalize vector spaces and abelian groups, but can have torsion, relations, and no basis.}{Commutative algebra, representation theory, homological algebra, algebraic geometry, topology, and differential geometry.}{Modules generalize vector spaces over rings, so generators, relations, maps, and exact sequences describe algebra beyond fields.}

\paragraph{Emmy Noether}
Noether developed the chain-condition methods that organize finitely generated modules and ideals. Her structural viewpoint made it possible to study algebra through generators, relations, and maps rather than through coordinates alone.

\paragraph{Emil Artin, Saunders Mac Lane, and homological algebraists}
The structure theorem over principal ideal domains classified important modules, while homological algebra organized modules through exact sequences, derived functors, and resolutions. These developments made modules central to representation theory, algebraic geometry, and topology \cite{module_theory_ref1}.
\end{historylist}
\section*{3. Theory}
\subsection*{Core theory}
\paragraph{The problem and definition}
Vector spaces allow addition and multiplication by scalars from a field. Many natural coefficient systems are rings instead: integers, polynomial rings, or smooth functions. A module over a ring $R$ is an additive abelian group $M$ with scalar multiplication satisfying distributivity, associativity with ring multiplication, and the identity rule. Every abelian group is a $\mathbb Z$-module, and every vector space is a module over its field \cite{module_theory_ref1}.

The difference is important. In a vector space, every nonzero scalar is invertible and bases behave uniformly. In a module, $2m$ may be zero for nonzero $m$, and a generating set may have relations.

\paragraph{Examples}
The ring $R$ is a module over itself, and $R^n$ is a free module. The integers $\mathbb Z/n\mathbb Z$ are $\mathbb Z$-modules with torsion. Ideals are submodules of a ring. Quotient rings are quotient modules.

If $X$ is a smooth manifold, smooth functions $C^\infty(X)$ form a ring. Smooth vector fields, differential forms, and tensor fields form modules over this ring. Sections of a vector bundle form projective modules; Swan's theorem says that over suitable compact spaces projective modules correspond to vector bundles.

\paragraph{Submodules, maps, and quotients}
A submodule is an additive subgroup closed under multiplication by ring elements. A module homomorphism preserves addition and scalar multiplication. Its kernel is a submodule, and the quotient by the kernel is isomorphic to its image.

The submodule generated by a set consists of finite sums $\sum r_i x_i$. A module is finitely generated when a finite set generates it, finitely presented when finitely many generators and relations describe it, and free when it has a basis. Projective modules behave like direct summands of free modules; injective modules provide extension properties.

\paragraph{Structure and relations}
Over a principal ideal domain, finitely generated modules split into a free part and cyclic torsion parts. The integer case is the structure theorem for finitely generated abelian groups:
\[
M\cong\mathbb Z^r\oplus\mathbb Z/d_1\oplus\cdots\oplus\mathbb Z/d_k,
\]
with divisibility conditions on the $d_i$. Over arbitrary rings this simple classification fails, so one uses exact sequences, resolutions, localization, and invariants such as rank and depth.

Modules over matrix rings are equivalent in category to modules over the original ring in the Morita theory. A Lie-algebra module is a module over its universal enveloping algebra. Representations of groups are modules over group algebras.

\paragraph{Homological algebra and geometry}
A projective or free resolution replaces a complicated module by a chain of simple modules. Applying $\operatorname{Hom}$ or tensor product to a resolution produces Ext and Tor, which measure failure of exactness. This is the algebra behind derived functors and sheaf cohomology.

In algebraic geometry, a sheaf of modules records functions or sections locally on open sets. Coherent modules over a coordinate ring describe algebraic vector bundles, ideals, and singularities. Localization studies a module near a prime ideal, while completion studies its formal neighborhood \cite{module_theory_ref2}.

\paragraph{What the theory depends on and where it is useful}
Module theory depends on rings, groups, linear algebra, and exact sequences. It helps representation theory turn actions into modules, algebraic geometry encode sheaves, topology calculate homology, and K-theory classify projective modules.

\section*{4. Important theorems and results}
\begin{enumerate}[leftmargin=*,itemsep=0.35em]
\item \textbf{Structure theorem over a PID.} Finitely generated modules decompose into free and cyclic torsion parts.

\item \textbf{First isomorphism theorem.} A module modulo the kernel of a homomorphism is isomorphic to its image.

\item \textbf{Nakayama's lemma.} Over a local ring, generators modulo the maximal ideal lift to generators of a finitely generated module.

\item \textbf{Hilbert syzygy theorem.} Polynomial rings over a field have finite projective resolutions of bounded length.

\item \textbf{Swan's theorem.} Suitable finitely generated projective modules over continuous functions are exactly vector-bundle sections \cite{module_theory_ref3}.
\end{enumerate}

\theoryapplications
\theoryconnections{module theory}{%
\item Linear algebra supplies vector spaces as the simplest modules, while ring theory supplies scalar multiplication and ideals.
\item Homological algebra studies exact sequences and resolutions, and algebraic geometry interprets modules as sheaves or sections.
}{%
\item Module theory helps classify algebraic objects, solve systems over rings, and compute invariants such as Ext and Tor.
\item It also provides the language for representations, coherent sheaves, and finitely generated data in algebraic geometry.
}

\section*{Glossary}
\begin{description}[leftmargin=*,style=nextline,itemsep=0.3em]
\item[\textbf{Module.}] An abelian group equipped with a scalar multiplication by elements of a ring, generalizing the idea of a vector space to scalars that may not be invertible.
\item[\textbf{Torsion.}] An element $m$ of a module is torsion if some nonzero ring element annihilates it (i.e., $rm=0$); torsion modules have elements that cannot be part of a basis.
\item[\textbf{Free module.}] A module that has a basis, meaning every element can be written uniquely as a finite linear combination of basis elements with ring coefficients.
\item[\textbf{Projective module.}] A module that is a direct summand of a free module; it behaves like a vector bundle in that it is locally free even if it is not globally free.
\item[\textbf{Exact sequence.}] A sequence of modules and homomorphisms in which the image of each map equals the kernel of the next, encoding precise relationships between algebraic objects.
\item[\textbf{Localization.}] The process of making selected ring elements invertible, producing a new ring and module focused on behavior near a chosen prime ideal or region of the spectrum.
\item[\textbf{Injective module.}] A module $I$ such that any homomorphism from a submodule into $I$ extends to the whole module.
\item[\textbf{Ext and Tor.}] Functors measuring the failure of $\operatorname{Hom}$ and tensor product to preserve exact sequences.
\item[\textbf{Principal ideal domain (PID).}] An integral domain in which every ideal is generated by a single element.
\item[\textbf{Coherent sheaf.}] A sheaf of modules locally generated by finitely many sections with finitely many relations.
\end{description}

\section*{7. Sample Questions}
\emph{Exercise reference:} \cite{module_theory_ref1}. The cited edition is the source for the exercise set; consult its Exercises section for the official statement and numbering.
\begin{enumerate}[leftmargin=*,itemsep=0.9em]
\item \textbf{Exercise 1.} Classify all finitely generated abelian groups of order $8$ up to isomorphism.

\emph{Solution.} By the structure theorem, the groups are: $\mathbb{Z}/8\mathbb{Z}$ (invariant factors $8$), $\mathbb{Z}/4\mathbb{Z}\oplus\mathbb{Z}/2\mathbb{Z}$ (invariant factors $4,2$), and $\mathbb{Z}/2\mathbb{Z}\oplus\mathbb{Z}/2\mathbb{Z}\oplus\mathbb{Z}/2\mathbb{Z}$ (invariant factors $2,2,2$). These are the only partitions of the prime factorization $2^3$: namely $3$, $2+1$, and $1+1+1$.

\item \textbf{Exercise 2.} Compute $\mathbb{Z}/2\mathbb{Z}\otimes_{\mathbb{Z}}\mathbb{Z}/3\mathbb{Z}$.

\emph{Solution.} In $\mathbb{Z}/2\mathbb{Z}\otimes\mathbb{Z}/3\mathbb{Z}$, for any $a\otimes b$ we have $a\otimes b=3(a\otimes b)/3=a\otimes(3b/3)$. But $3b\equiv 0\pmod{3}$, so $a\otimes b=a\otimes 0=0$. Hence the tensor product is $\{0\}$, which is isomorphic to $\mathbb{Z}/\gcd(2,3)\mathbb{Z}=\mathbb{Z}/1\mathbb{Z}=0$.

\item \textbf{Exercise 3.} Find $\operatorname{Ann}_{\mathbb{Z}}(\bar{1})$ in $\mathbb{Z}/12\mathbb{Z}$, and determine all group homomorphisms $\varphi\colon\mathbb{Z}/12\mathbb{Z}\to\mathbb{Z}/8\mathbb{Z}$.

\emph{Solution.} $\operatorname{Ann}(\bar{1})=\{n\in\mathbb{Z}:n\cdot\bar{1}=\bar{0}\}=12\mathbb{Z}$. A homomorphism $\varphi$ is determined by $\varphi(\bar{1})=\bar{g}\in\mathbb{Z}/8\mathbb{Z}$ with $12\bar{g}=\bar{0}$, i.e., $8\mid 12g$, i.e., $2\mid g$. The solutions are $g\in\{0,2,4,6\}$, giving $4$ homomorphisms.

\item \textbf{Exercise 4.} Compute $\operatorname{Ext}^1_{\mathbb{Z}}(\mathbb{Z}/2\mathbb{Z},\,\mathbb{Z}/4\mathbb{Z})$ using the free resolution $0\to\mathbb{Z}\xrightarrow{\times 2}\mathbb{Z}\to\mathbb{Z}/2\mathbb{Z}\to 0$.

\emph{Solution.} Apply $\operatorname{Hom}(-,\mathbb{Z}/4\mathbb{Z})$: $\operatorname{Hom}(\mathbb{Z},\mathbb{Z}/4\mathbb{Z})\cong\mathbb{Z}/4\mathbb{Z}$ and the induced map is $\times 2\colon\mathbb{Z}/4\mathbb{Z}\to\mathbb{Z}/4\mathbb{Z}$. The image is $\{0,2\}\cong\mathbb{Z}/2\mathbb{Z}$. $\operatorname{Ext}^1$ is the cokernel: $\mathbb{Z}/4\mathbb{Z}/\{0,2\}\cong\mathbb{Z}/2\mathbb{Z}$.

\item \textbf{Exercise 5.} Compute $\operatorname{Tor}_1^{\mathbb{Z}}(\mathbb{Z}/2\mathbb{Z},\,\mathbb{Z}/4\mathbb{Z})$ using the same free resolution.

\emph{Solution.} Tensor the resolution with $\mathbb{Z}/4\mathbb{Z}$: $\mathbb{Z}/4\mathbb{Z}\xrightarrow{\times 2}\mathbb{Z}/4\mathbb{Z}$. The kernel of $\times 2$ on $\mathbb{Z}/4\mathbb{Z}$ is $\{0,2\}\cong\mathbb{Z}/2\mathbb{Z}$. Hence $\operatorname{Tor}_1\cong\mathbb{Z}/2\mathbb{Z}$.

\end{enumerate}
\section*{8. References}
\addcontentsline{toc}{section}{References}

\begin{thebibliography}{9}
\bibitem{module_theory_ref1} David Eisenbud, \emph{Commutative Algebra with a View Toward Algebraic Geometry}, Springer, 1995.
\bibitem{module_theory_ref2} Joseph J. Rotman, \emph{An Introduction to Homological Algebra}, 2nd ed., Springer, 2009.
\bibitem{module_theory_ref3} Charles A. Weibel, \emph{An Introduction to Homological Algebra}, Cambridge University Press, 1994.
\end{thebibliography}